\section{Real monomial data in parameter families}\label{sec:real}
We work first with a compact semialgebraic model $X$, a continuous
semialgebraic probability map $F$, and a continuous semialgebraic monic
target polynomial $A_x$ of fixed degree $n$. Complex coefficients are
represented by real and imaginary parts. Suppose the target is constant
on $F^{-1}(P)$, with value $A_0(z)=\prod_r(z-r)^{m_r}$. Disjoint
isolating discs and compactness give, on a sufficiently small closed
observation neighbourhood, unique monic cluster factors
\begin{equation}\label{eq:clusters}
 A_{r,x}(r+z)=z^{m_r}+\sum_{a<m_r}c_{r,a}(x)z^a,
 \qquad k_{r,a}=m_r-a.
\end{equation}
Multiplication of coprime monic factors has invertible differential.
The cluster factors are consequently analytic functions of the total
coefficients. Their graphs with the isolating conditions are semialgebraic.

\begin{definition}[Input and families]\label{def:input}
An algebraic input is a finite first-order polynomial presentation of
$X,F,A,P$, design weights, isolating data, denominator exclusions and
branch choices. A real algebraic constant is encoded by an integer
polynomial and a rational isolating interval; a complex constant is
encoded by two real algebraic constants. A fixed-format family has a
semialgebraic parameter set $\Theta$, a semialgebraic
$X\subset\Theta\times K$ with $K$ compact and compact fibres, continuous
semialgebraic maps $F,A$ on $X$, and a semialgebraic choice
$P_\theta\in F_\theta(X_\theta)$. Degrees and numbers of coordinates
are fixed. Effectivity refers to algebraic presentations and numerical
outputs to algebraic parameter values. No complexity bound is asserted.
\end{definition}

Adjoin nonnegative square-root coordinates to the observation graph and
write
\[
 R_{je}=\sqrt{w_j}(\sqrt{F(x)_{je}}-\sqrt{P_{je}}),\quad
 G=\sum_{j,e}R_{je}^2,\quad
 H_{r,a}=(\Re c_{r,a})^2+(\Im c_{r,a})^2.
\]
They are polynomial functions on the resulting compact real graph
$\Gamma$. Zero probability cells are permitted. No passage to a Fisher
quadratic form is used in this construction.

\begin{definition}[Real monomial presentation]\label{def:monomial}
A real monomial presentation is a finite collection of smooth real
algebraic sources and proper maps, collectively covering a neighbourhood
of $G=0$ in $\Gamma$, with finite algebraic etale coordinate
presentations. On each source component either $G\equiv0$ and all
$H\equiv0$, or
\begin{equation}\label{eq:monomial}
 G=u\prod_i|y_i|^{a_i},\qquad
 H_{r,a}\equiv0\quad\hbox{or}\quad
 H_{r,a}=v_{r,a}\prod_i|y_i|^{b_{r,a,i}}.
\end{equation}
The units have positive upper and lower bounds on smaller closed real
coordinate boxes above $G=0$. A contributing divisor has $a_i>0$ and
a real point away from the other divisors, with an admissible transversal
on which $G>0$. An identically zero $H$ has order $+\infty$; a component
$G\equiv0$ is recorded separately and contributes no ratio. Properness
belongs to the source maps, not to an open coordinate box. The finite
encoding of an adapted coordinate presentation is specified in
Appendix~\ref{app:foundations}.
\end{definition}

\begin{lemma}[Spreading a finite resolution certificate]\label{lem:spread}
Let $S=\operatorname{Spec}A$ be an integral affine variety over a
characteristic-zero field, let $K=\operatorname{Frac}(A)$, and let
$Y\to S$ be of finite type. Suppose that over $K$ there is a finite
reduced closed chain
\[
 Y_{0,K}=Y_{K,\mathrm{red}}\supset Y_{1,K}\supset\cdots
                 \supset Y_{q,K}=\varnothing,
\]
projective maps $Z_{j,K}\to Y_{j,K}$ with smooth sources, specified
inverse isomorphisms over $Y_{j,K}\setminus Y_{j+1,K}$, and finite
smooth-centre principalizations of specified functions on those sources.
There is a nonempty open $S^\circ$ on which the following hold on every
fibre. The closed chain is flat with geometrically reduced fibres and
ends in the empty set; the inverse maps remain inverse over its actual
successive complements. Sources and centres are relatively smooth,
principalization blow-ups commute with fibre formation, and the final
sources have relative normal-crossing divisors with fixed integer
monomial orders on a finite adapted etale atlas. Zero identities and
inverse units persist. A nonzero monomial cannot become identically zero
on a component of a retained fibre. For algebraic input, such finite
data and an exceptional set are effectively constructible.
\end{lemma}
\begin{proof}
Lemma~\ref{lem:finite-diagram} descends the complete diagram, including
the actual open complements and both inverse identities, and imposes
flatness and geometric reducedness. Lemma~\ref{lem:rees} proves the
needed base change for all powers of every centre ideal after one
localization per finite collection of charts. The explicit bad-locus
construction in Appendix~\ref{app:foundations} treats relative
smoothness, normal-crossing intersections, chart coverage, units and
zero flags separately, and proves that each discarded algebraic image
has proper closure. The nonzero-specialization assertion follows from
relative codimension one and the inverse units, not from generic
nonvanishing alone. The constructive certificate specification in
Appendix~\ref{app:effective} proves the last assertion.
\end{proof}

\begin{theorem}[Relative real principalization]\label{lem:relative}
On the exact-target identifying locus of a fixed-format family there is
a finite semialgebraic stratification with fixed cluster multiplicities,
a fixed finite collection of algebraic source chart presentations, fixed
integer orders and fixed real-accessibility flags as in
Definition~\ref{def:monomial}. The source maps are fibrewise proper real
covers. Charts, units and box sizes may vary with the parameter. On
compact subsets of a stratum the cover is usable over a common
sufficiently small observation neighbourhood. The construction is
effective for the algebraic presentations of Definition~\ref{def:input}.
\end{theorem}
\begin{proof}
Lemma~\ref{lem:slack-exact} algebraizes the proper enlarged graph by
square-slack lifts without discarding strict exclusions or real branches.
Resolve the reduced generic fibre and its successive reduced singular
loci, then principalize the nonzero functions on the smooth sources.
Lemma~\ref{lem:spread} spreads this finite certificate. The explicit
exceptional algebraic loci in Appendix~\ref{app:foundations} are
reconstructed on reduced bases of strictly smaller dimension, giving
finite termination. Lemma~\ref{lem:real-descent} proves both inclusions
in the exact real coverage assertion, including lower-dimensional real
loci, and descends the selected label branches. Lemma
\ref{lem:real-accessibility} supplies the real transversals and the finite
real flag partition. That partition is not confused with the algebraic
dimension-drop recursion. The final proof in
Appendix~\ref{app:foundations} supplies the uniform neighbourhood and
unit bounds by properness over compact parameter subsets. These
arguments prove the asserted cover on the original graph, rather than
on an algebraic relaxation of its inequalities.
\end{proof}

\begin{proposition}[Fixed-centre real cover]\label{prop:real-cover}
Every graph $\Gamma$ above has a real monomial presentation. Exact
pieces and lower-dimensional real loci are included. Every recorded
divisor has a real admissible attaining arc.
\end{proposition}
\begin{proof}
Apply Theorem~\ref{lem:relative} to a one-point base and use
Lemma~\ref{lem:real-accessibility}.
\end{proof}

\begin{corollary}[Intrinsic scalar orders]\label{cor:orders}
For each nonrigid cluster coefficient put
\begin{equation}\label{eq:exponent}
 \beta_{r,a}=\min_{\text{recorded real divisors}}\frac{b_{r,a,i}}{a_i},
 \qquad \alpha=\min_{r,a}\frac{\beta_{r,a}}{k_{r,a}}.
\end{equation}
Either the target is constant on an observation neighbourhood, or
$\alpha>0$ is rational and $\omega_P(t)\asymp t^\alpha$.
Every finite coefficient minimum is attained by an admissible arc.
The numbers are independent of the monomial presentation.
\end{corollary}
\begin{proof}
The complete real divisor list and the monomial inequalities in
Lemma~\ref{lem:real-accessibility} give $H\le C G^\beta$, hence
$|c_{r,a}|\le C h_w^\beta$. A minimizing transversal satisfies
$G\asymp u^{a_i}$ and $H\asymp u^{b_i}$ and attains the exponent.
The inclusion $\{G=0\}\subseteq\{H=0\}$ forces positive contributing
orders. If all coefficients vanish on an observation neighbourhood, the
target is locally rigid. Otherwise, for a monic polynomial of degree $m$,
its greatest root modulus is comparable to
$\max_{a<m}|c_a|^{1/(m-a)}$. Elementary symmetric identities give one
bound; dominance of $z^m$ outside a fixed multiple of this maximum gives
the other. The centre lies in every ball, and the diameter is between
its greatest centre displacement and twice that displacement. This
proves the order of $\omega_P$. The optimal coefficient inequality is
intrinsic, so its attained exponent is independent of the presentation.
The scalar divisor-ratio calculation is classical; compare \cite{BE09}.
\end{proof}

\section{Joint specialization and uniform spectral asymptotics}\label{sec:joint}
On a nonrigid identified stratum write $\alpha=p/q$ in lowest terms.
Let $\Psi(\theta,\varepsilon,x,U,V)$ denote graph membership and
\begin{equation}\label{eq:weighted}
 R_{je}(\theta,x)=\varepsilon^q U_{je},\qquad
 c_{r,a}(\theta,x)=\varepsilon^{p k_{r,a}}V_{r,a},\qquad
 \sum_{j,e}U_{je}^2\le1.
\end{equation}
Complex equalities denote two real equalities. Define the joint family
by the fixed-centre formula
\begin{align}
 (\theta,U,V)\in\mathcal J\quad\Longleftrightarrow\quad
 &\theta\in S\ \text{and}\ \forall\rho>0\ \exists\varepsilon,x,U',V':
 \nonumber\\
 &0<\varepsilon<\rho,\quad\Psi(\theta,\varepsilon,x,U',V'),\quad
 \|(U',V')-(U,V)\|^2<\rho^2.\label{eq:closure-formula}
\end{align}
Put $\cI_0(\theta)=\mathcal J_\theta$ and
$\cC_0(\theta)=\pr_V\cI_0(\theta)$. There is no quantified
$\theta'$ in this formula. Closure with a drifting centre is a different
operation, and taking independent closures in the coefficient coordinates
is also a different operation.

\begin{theorem}[The definable joint initial family]\label{thm:joint}
The sets $\cI_0(\theta)$ and $\cC_0(\theta)$ form semialgebraic families
of nonempty compact sets. The coefficient images of the radius-$t$
observation ball, scaled by $c_{r,a}/t^{\alpha k_{r,a}}$, converge in
Hausdorff distance to $\cC_0(\theta)$ for every fixed $\theta$. If
\[
 Z_r(V)=Z\left(z^{m_r}+\sum_{a<m_r}V_{r,a}z^a\right),\qquad
 \mathcal L_0(\theta)=\{(Z_r(V))_r:V\in\cC_0(\theta)\},
\]
then the scaled cluster-root images converge to $\mathcal L_0(\theta)$,
and
\begin{equation}\label{eq:joint-diameter}
 \omega_\theta(t)=C(\theta)t^\alpha+o(t^\alpha),\qquad
 C(\theta)=\max_{V,V'\in\cC_0(\theta)}\max_r
 d_\infty(Z_r(V),Z_r(V'))\in(0,\infty).
\end{equation}
This is the diameter of a joint feasible image, not that of a Cartesian
product of scalar envelopes.
\end{theorem}
\begin{proof}
Corollary~\ref{cor:orders} bounds every scaled coefficient for small $t$;
$U$ is already bounded. The fixed-centre closure is nonempty and compact.
Quantifier elimination proves semialgebraicity with $\theta$ as a free
parameter. Compactness of the graph fibre lets one pass to convergent
subsequences of the unscaled graph variables when projecting an initial
point, so projection loses no limiting coefficient point.

Fix $\theta$ and a point of the outer coefficient limit. Its distance
to the positive-radius coefficient section is a bounded nonnegative
semialgebraic function of $t$, with a subsequence tending to zero. The
one-variable monotonicity theorem makes its limit zero. A finite net of
the compact outer limit gives the inner Hausdorff inclusion. Failure of
the outer inclusion would give a bounded convergent subsequence of scaled
coefficients outside the defining limit, a contradiction. This proves
convergence of the entire ball. Continuity of monic root multisets on
bounded coefficient sets gives the corresponding root-set convergence.

Let $\Delta$ be the least separation of distinct base roots. For small
$t$, each cluster lies in its $\Delta/4$ disc. A matching within clusters
has error tending to zero; a cross-cluster matching has an edge of length
at least $\Delta/2$. An optimal matching therefore stays within clusters.
Translation and scaling there are exact. Taking the diameter gives
\eqref{eq:joint-diameter}. An attaining coefficient transversal prevents
the limiting diameter from being zero. This also explains why all
inequality faces and joint relations in the closure remain relevant.
\end{proof}

\begin{proposition}[Semialgebraic decay on strata]\label{prop:decay}
Let $e(\theta,t)\ge0$ be semialgebraic for $0<t<\tau(\theta)$, with
$\tau>0$ semialgebraic, and suppose $e(\theta,t)\to0$ for each fixed
$\theta$. After a finite Nash stratification, on every stratum a rational
$\nu>0$ and a positive continuous semialgebraic $\tau_0(\theta)$ satisfy
\[
 0<t<\tau_0(\theta)\quad\Longrightarrow\quad e(\theta,t)\le t^\nu.
\]
The estimate is uniform on compact subsets; $\nu$ need not be sharp.
\end{proposition}
\begin{proof}
Lemma~\ref{lem:support-order} gives a finite positive rational candidate
list from the polynomial supports of a quantifier-free graph formula,
including specialized zero coefficients and the eventual-zero case.
Choose $\nu$ below its positive candidates. The complete argument in
Appendix~\ref{app:effective} selects a genuinely positive semialgebraic
threshold, makes it continuous, and takes its minimum only on a compact
subset of a stratum. No uniform rate is inferred from pointwise
convergence alone.
\end{proof}

\begin{theorem}[Uniform spectral power law]\label{thm:uniform}
The identified nonrigid locus has a finite Nash stratification on which
$\alpha$ is fixed, $\mathcal L_0$ has constant semialgebraic topological
type, and $C(\theta)$ is positive and Nash. For every compact subset $H$
of a stratum, for some $\nu>0$,
\begin{equation}\label{eq:uniform-law}
 d_H(\mathcal L_t(\theta),\mathcal L_0(\theta))=O_H(t^\nu),\qquad
 \omega_\theta(t)=C(\theta)t^\alpha+O_H(t^{\alpha+\nu}).
\end{equation}
Here $\mathcal L_t$ is the exact clustered root image scaled by
$t^{-\alpha}$ and $d_H$ uses the maximum cluster bottleneck metric.
\end{theorem}
\begin{proof}
Represent each root multiset by all its labelled arrays through elementary
symmetric identities. The formula
\begin{equation}\label{eq:matching-formula}
 M_d(z,z')=\bigvee_{\sigma\in\mathfrak S_m}
              \bigwedge_{i=1}^m |z_i-z'_{\sigma(i)}|^2\le d^2
\end{equation}
expresses bottleneck distance at most $d$. Use the conjunction over
clusters, and real and imaginary parts for complex roots. The two
inclusions
\[
 (\forall z\in\mathcal L_t)(\exists z'\in\mathcal L_0)M_d(z,z'),
 \qquad
 (\forall z'\in\mathcal L_0)(\exists z\in\mathcal L_t)M_d(z,z')
\]
and minimality over $d\ge0$ define the Hausdorff error. Compactness gives
attainment. Rational powers of $t$ are encoded by positive algebraic
roots, so this error is semialgebraic. Theorem~\ref{thm:joint} proves
its pointwise limit zero; Proposition~\ref{prop:decay} gives the
compact-uniform power bound.

The diameter has a first-order formula saying that all pairs have distance
at most $C$ and one pair attains it. It is therefore semialgebraic and
positive. Lemma~\ref{lem:nash-graph} gives its Nash refinement. Refine
also by positive cluster separation and observation-radius choices, so
that compact minima exclude cross-cluster matching uniformly. Hardt
triviality supplies the topological types separately. Finally diameters
differ by at most twice Hausdorff distance, which gives the second
estimate. Topological triviality has supplied neither the error rate nor
the Nash dependence.
\end{proof}

\begin{theorem}[A finite real spectral atlas]\label{thm:atlas}
A fixed-format family admits a finite semialgebraic stratification
retaining the following data: the semialgebraic type and dimension of the
exact target coefficient fibre and its root image; on the identified
locus, rigidity or a positive rational exponent, cluster multiplicities,
and a fixed real monomial order list with accessibility and minimizing
indices; on its nonrigid part, the joint initial coefficient and root-set
types, a Nash leading diameter and the compact-uniform power remainder;
and finite algebraic descriptions of these objects. On a subfamily with
a proved Fisher-cone representation, the active supports and choices of
leading constants can be retained in the same refinement. On a
nonidentified stratum the modulus tends to the positive diameter of the
exact target root image. No general leading set is asserted to be
quadratic, and no local leading constant excludes a remote exact component.
\end{theorem}
\begin{proof}
The exact coefficient fibre is $\{A_x:F_\theta(x)=P_\theta\}$, not the
parameter fibre itself. Apply Hardt triviality to its compact image and
root graph. Singleton images are selected by a first-order formula.
At fixed $\theta$, compactness shows that decreasing observation balls
converge in target image to the exact image, so their diameters converge
to its diameter. On the singleton locus apply Theorem~\ref{lem:relative},
Corollary~\ref{cor:orders}, and Theorems~\ref{thm:joint} and
\ref{thm:uniform}. These are separate constructions entering one finite
refinement. A proved Fisher subfamily contributes only finitely many
additional support and comparison conditions. Effective descriptions are
provided by the following theorem and Appendix~\ref{app:effective}.
\end{proof}

\section{Effective output and the role of computation}\label{sec:effective}
\begin{theorem}[Effective algebraic specialization]\label{thm:effective}
For an algebraic input on the identified locus, finite algorithms compute
local rigidity, the nonrigid exponent $\alpha$, quantifier-free
presentations of $\cI_0,\cC_0,\mathcal L_0$, and an integer polynomial
with a rational isolating interval for $C$. Exact identification itself
is decidable. In algebraically presented parameter families the
corresponding strata and formulas are effectively describable. These
assertions are independent of any finite example regression.
\end{theorem}
\begin{proof}
Appendix~\ref{app:effective} gives the complete input--output procedure.
It constructs each maximum $m_H(s)=\max\{H:G\le s\}$ by elimination,
extracts a finite support-pair list, and tests the comparability formula
\eqref{eq:order-decision} for every candidate. The unique true candidate
selects the actual positive maximum-function branch, while a separate
formula detects eventual zero. No unspecified branch-isolation oracle is
used. Substitute the resulting $p/q$ in \eqref{eq:closure-formula},
eliminate, adjoin roots and use \eqref{eq:matching-formula} to compute the
leading diameter. Real-root isolation and field-norm conversion provide
its integer polynomial and isolating interval. With free parameters the
same finite tests compute strata, followed by Nash refinement. A second,
explicit finite-certificate search for the geometric presentations is
given in that appendix; spectral outputs do not depend on executing it.
\end{proof}

The symbolic calculations accompanying the explicit polynomial examples
test identities and specified finite optimization problems. They do not
implement the universal algorithms above and are not proof certificates
for relative principalization, uniform asymptotics, or arbitrary degree.
